\section{Where ``starts closer'' is true: parameter space}
\label{sec:param}

\Cref{thm:marginal} shows that a larger random model does not start with a
smaller loss. It can still start nearer to a fit, if nearness is measured in the
space of parameters. For models that are linear in their parameters this is
elementary, and it is the lazy-training picture of
\citet{jacot2018ntk,chizat2019lazy} in its simplest case. We state it because it
is the form in which the original idea holds.

\subsection{Setting}

Let $\Theta$ and $Y$ be real Hilbert spaces and $A:\Theta\to Y$ a bounded linear
map. We read $\theta\in\Theta$ as the parameters, $A\theta$ as the outputs of
the model on the data, and $y\in Y$ as the labels. The \emph{residual} is
$r(\theta)=y-A\theta$, a \emph{fit} is a $\theta$ with $A\theta=y$, and the
\emph{travel} from $\theta_0$ is the distance to the nearest fit,
\[
  \tau(\theta_0)=\inf\{\|\theta-\theta_0\|:A\theta=y\}.
\]
A Fourier series is such a model in its coefficients, and so is a neural
network in the weights of its last layer when the earlier layers are held
fixed. The \emph{kernel} is $K=AA^*$.

\subsection{Results}

\begin{proposition}[One loss curve]\label[proposition]{prop:path}
If $A\theta_1=y$, then $r\bigl(\theta_0+t(\theta_1-\theta_0)\bigr)=(1-t)\,r(\theta_0)$
for every $t\in\R$. Consequently, if $d=\|\theta_1-\theta_0\|>0$ and
$0\le\varepsilon\le d$, there is a $\theta$ with $\|\theta-\theta_0\|=\varepsilon$
and $\|r(\theta)\|=(1-\varepsilon/d)\,\|r(\theta_0)\|$.
\end{proposition}

\begin{proof}
$y-A\theta_0-t(A\theta_1-A\theta_0)=(1-t)(y-A\theta_0)$. For the second
statement take $t=\varepsilon/d$.
\end{proof}

Every model walks the same loss curve from its start to a fit; what differs is
the length of the walk. A fit at distance $d$ means that a movement of
$\varepsilon$ removes the fraction $\varepsilon/d$ of the residual.

\begin{theorem}[A model is its kernel]\label{thm:kernel}
\emph{(a)} One step of gradient descent on $\tfrac12\|y-A\theta\|^2$ with step
size $\eta$, $\theta^+=\theta+\eta A^*r(\theta)$, gives
$r(\theta^+)=(I-\eta K)\,r(\theta)$.
\emph{(b)} If $K$ is invertible, then
$\tau(\theta_0)=\|A^*K^{-1}r(\theta_0)\|=\sqrt{\langle K^{-1}r(\theta_0),r(\theta_0)\rangle}$.
Hence two models with the same kernel and the same initial residual have the
same residual after every gradient step, and the same travel.
\end{theorem}

\begin{proof}
(a) $y-A\theta-\eta AA^*r(\theta)=r(\theta)-\eta Kr(\theta)$.
(b) Write $r_0=r(\theta_0)$ and $\delta^\ast=A^*K^{-1}r_0$. Then
$A\delta^\ast=r_0$, so $\theta_0+\delta^\ast$ is a fit. If $\theta_0+\delta$ is
any fit, then $A\delta=r_0$ and
\[
  \|\delta^\ast\|^2=\langle K^{-1}r_0,AA^*K^{-1}r_0\rangle
  =\langle K^{-1}r_0,r_0\rangle=\langle K^{-1}r_0,A\delta\rangle
  =\langle\delta^\ast,\delta\rangle\le\|\delta^\ast\|\,\|\delta\|,
\]
so $\|\delta^\ast\|\le\|\delta\|$. The second and third expressions show that
$\|\delta^\ast\|^2=\langle K^{-1}r_0,r_0\rangle$, which involves only $K$ and
$r_0$.
\end{proof}

\begin{corollary}[The exchange rate]\label[corollary]{cor:rate}
If $K=c\,I$ with $c>0$, that is, if the features are orthogonal on the data side
with squared length $c$, then
\[
  \tau(\theta_0)=\frac{\|r(\theta_0)\|}{\sqrt c},
\]
the residual after $t$ gradient steps is $(1-\eta c)^t\,r(\theta_0)$, and with
$\eta=1/c$ a single step reaches a fit. For $c=1$ the travel equals the size of
the residual.
\end{corollary}

Loss and travel are then one quantity in two units, and the number of
parameters enters only through the rate $\sqrt c$ between them.

\begin{proposition}[More parameters never move the fit away]\label[proposition]{prop:more}
Let $B:\Theta'\to Y$ be a second bounded linear map and consider the larger
model $(\theta,\theta')\mapsto A\theta+B\theta'$. If the smaller model has a
fit, the travel of the larger model from $(\theta_0,0)$ is at most the travel
of the smaller model from $\theta_0$.
\end{proposition}

\begin{proof}
If $A\theta=y$ then $(\theta,0)$ is a fit of the larger model at the same
distance from $(\theta_0,0)$.
\end{proof}

Together with \cref{thm:marginal} this gives the corrected form of the
conjecture, for models linear in their parameters: with no movement allowed, a
larger random model is no closer to the data; with any positive amount of
movement allowed, it is at least as close. \Cref{prop:more} starts the new
parameters at zero. The experiment below starts them at random.

\subsection{A family in which the number of parameters is the only dial}
\label{sec:fourier}

The sine family has one random number, so it cannot test a statement about the
number of parameters. We therefore take a Fourier series with $m$ terms,
\[
  f_\theta(x)=a_0+\sum_{k=1}^{m}\bigl(a_k\cos(k\pi x)+b_k\sin(k\pi x)\bigr),
\]
with $P=2m+1$ parameters drawn independently from $N(0,V/(m+1))$, $V=1/3$. The
output then has variance $V$ at every $x$ and every $m$, so the one-point
marginals do not depend on $m$. The data are $n=20$ points with $x$ and $y$
uniform on $[-1,1]$; a fit exists once $P\ge n$. For each $m$ we use 20 data
sets and 20 random starts, and measure the initial squared error, the
nearest-point distance, the travel $\tau$, and the number of gradient steps
needed to reach a mean squared error of $10^{-3}$ with the step size
$1/\lambda_{\max}(K)$.

\begin{table}[t]
  \caption{A random Fourier series with $P$ parameters on 20 points with random
  labels. Initial squared error and nearest-point distance are means; travel
  and steps are medians. The last column is $\sqrt{n\cdot\mathrm{MSE}/(m+1)}$,
  the travel that \cref{cor:rate} gives if the features were exactly orthogonal
  with $c=m+1$. The nearest-point distance was not measured above $P=129$.}
  \label{tab:fourier}
  \centering
  \begin{tabular}{rccccc}
    \toprule
    $P$ & initial MSE & nearest-point & travel & gradient steps & $\sqrt{n\,\mathrm{MSE}/(m+1)}$ \\
    \midrule
      9 & 0.65 & 0.229 & no fit & never & -- \\
     17 & 0.71 & 0.157 & no fit & never & -- \\
     21 & 0.65 & 0.128 & 448    & $1.9\times10^{8}$ & -- \\
     25 & 0.64 & 0.110 & 35.5   & $7.9\times10^{4}$ & -- \\
     33 & 0.66 & 0.090 & 16.0   & $3.3\times10^{4}$ & 0.88 \\
     65 & 0.71 & 0.056 & 1.61   & 166 & 0.66 \\
    129 & 0.66 & 0.027 & 0.48   & 13  & 0.45 \\
    257 & 0.66 & --    & 0.34   & 3   & 0.32 \\
    513 & 0.63 & --    & 0.23   & 3   & 0.22 \\
    \bottomrule
  \end{tabular}
\end{table}

\Cref{tab:fourier} shows three things. The initial squared error does not move:
it scatters around $\E y^2+V=2/3$, as \cref{thm:marginal} requires. The travel
and the number of steps fall steeply: from 21 to 513 parameters the travel falls
by a factor of about 2000, and training goes from $1.9\times10^{8}$ steps to 3.
And the worst case is at the threshold $P\approx n$, where a fit exists but is
far away; this is where double descent has its peak
\citep{belkin2019reconciling,nakkiran2020double}.

Far above the threshold the measured travel agrees with \cref{cor:rate}. The
kernel of these features has diagonal entries $m+1$ and, for scattered points
and $m\gg n$, comparatively small off-diagonal entries, so it is close to
$(m+1)I$. The last column of \cref{tab:fourier} is within $7\%$ of the measured
travel from $P=129$ upwards, and travel falls like $1/\sqrt P$. Near the
threshold the features are far from orthogonal and the formula fails by a wide
margin. That the kernel is close to $(m+1)I$ is an observation here, not a
proved statement.

The nearest-point distance also falls, roughly like $1/P$. This is
\cref{thm:sandwich} again: a random series with more terms of this scaling is
steeper. So in this family claims (G) and (P) both hold, and (L) still fails.
A larger random model is not nearly right. It is nearly adjustable.
